\documentclass[11pt]{article}
\usepackage{preprint}

\title{The Best Regime and the Worst Regime: A Further Review of Aristotle's Politics}
\author{Jiadi He\email{jh2288@georgetown.edu}}
\affil{Georgetown University}
\shorttitle{The Best Regime and the Worst Regime}
\keywords{political theory; regime; political community}

\begin{document}
\maketitle
\keywordsline

Aristotle is among the greatest scholars of humanity to date, with his work covering almost all disciplines and making great contributions to them. Politics is one of Aristotle's numerous works, but it is one of his most important. What is the best regime and what is the worst regime? This is one of the fundamental issues that Aristotle was most concerned about. His research on the city-state focused on what form is the best and how to achieve the ideal life that people imagine. It has long been a popular view that the best form of government in Aristotle's mind is monarchy. However, in recent studies of Aristotle's ideology, more and more scholars have come to believe that republican democracy is the ideal political system he pursues. This analysis will combine the latest research results on Aristotle in the international academic community to re-examine his political system theory, in order to determine what Aristotle's best and worst political systems are.

\section{The Political Community}

Aristotle believed that a city-state is a social organization of humanity, or a political community, which he called ‘\textit{polis},’ formed by people in pursuit of common interests (Aristotle, 1885). The primary characteristic of a city-state is its political nature, which distinguishes it from other social groups. As a political community, a city-state has a supreme nature among all social groups. He made a clear point in his book, \textit{Politics}:

\begin{quote}
“Every state is a community of some kind, and every community is established with a view to some good; for mankind always acts in order to obtain that which they think good. But, if all communities aim at some good, the state or political community, which is the highest of all, and which embraces all the rest, aims at good in a greater degree than any other, and at the highest good.” (Aristotle, 1885, p.3)
\end{quote}

The ‘\textit{polis’} referred to by Aristotle is neither a general city nor a modern nation state, but an early form of state. It already possesses the basic characteristics of a country: a certain range of fields, a certain scale of population, independent political institutions, and legitimate armed forces, among other features (Aristotle, 1885). Aristotle used this view to define what a political regime would be and the role politics played in creating a \textit{polis.}

Unlike modern political thinkers, Aristotle insists that city-states are not political societies deliberately created by humans, but rather human beings as a natural result of being a born political animal. The logic used by Aristotle to demonstrate the naturalness of a city-state is that in order for humans to reproduce, there is a natural union of men and women; Male and female spouses form a family, which is the basic form of society established by humans to meet natural needs; Several families form a village, and the village is a natural settlement derived from family reproduction; Several villages unite to form a city-state, which is a “self-sufficient” human organization and the highest and widest political community among all social groups. Within the city-state, some people are born rulers, while others are born ruled, which is also natural. Aristotle emphasized that all aspects, from human reproduction to the form of families and villages, to the emergence of city-states, as well as the formation of rulers and subjects, are the result of natural evolution. Therefore, his conclusion is very clear: “State is a creation of nature,” and the emergence of cities indicates that society evolves to a higher and more complete state (Aristotle, 1885).

Since a city-state is the pinnacle of human natural evolution and the highest completed form of human organization, it naturally has priority over any other form of human organization. More specifically, Aristotle concluded that city-states have unique priority in two senses. Firstly, compared to other forms of social organization such as families and villages, only the city-state is a self-sufficient and complete form of organization, and no other social group is complete. The city-state surpasses any other social group. Secondly, compared to individuals, a city-state is a whole, and individuals are parts. The whole must be higher than the parts, and parts only have meaning in the whole. (Young, 2000)

It is said that in order to systematically study the rise and fall history of city states and search for the best form of city states, Aristotle once compared and analyzed 158 ancient Greek city states. The current “Athenian Politics” is only one of the case studies of its 158 city-states. Therefore, Aristotle's city-state mainly refers to the form of social political organization that spread throughout ancient Greece. It includes a certain number of citizens, a certain scale of territory, and specific political governance methods. Aristotle referred to the political governance of city-states as the system of government, and believed that the system of government was the decisive factor of city-states and a substantive symbol of their mutual differences.

The original Greek word for ‘regime' is' politeiai ', while the English word is usually translated as' regime', ‘constitution', or ‘polity'. In Aristotle's view, a politeia is the arrangement of magistracies in a city-state, especially when it consists the highest of all. He also stated that the government is sovereign in a state, and the constitution is the government (Robinson, 1971). In the language of modern political science, a political system mainly refers to the organizational form and operational mechanism of a political power, including the nature and mode of production of the supreme ruler, the composition and operational mode of the political power. Qualitatively speaking, “political system” refers to the political system of a city-state (Robinson, 1971).

Aristotle further pointed out that there are two basic ways of political rule, one is to prioritize the interests of the rulers themselves, and the other is to prioritize the public interests of the city-state. He mentioned that there are two fundamentally different ways of governance: one is centered around the interests of the ruler, and the other is based on the interests of the ruled. The former is called ‘authoritarian rule' (master slave rule), and the latter is called ‘free man rule'. If all three forms of government can be based on the public interests of the city-state, then they are all conventional good governance. On the contrary, if they are not focused on pursuing the public interests of the city-state, but rather on pursuing the personal interests of the rulers themselves, then these three regimes will undergo distortion and tend towards their own opposite, becoming abnormal regimes (Robinson, 1971).

In other words, as Aquinas said, Aristotle actually divided all regimes into six types, of which monarchy, aristocracy and republic are conventional regimes, while tyrants, oligarchs and plebeians are variant regimes. Aquinas explains Aristotle's political system in this way: There are six different city-state organizations, where one or a few people or many people rule each city-state. If ruled by one person, it is either a monarchy or a tyranny. If he can achieve virtue and aim for the common interests of his subjects, it is a monarchy. If he is evil and causes everything to satisfy his own interests and harm the interests of his subjects, it is tyranny. If a few people should rule the city-state, and the person who takes care of the people's welfare is elected because of his virtue, we call it aristocracy system, or if a few people rule, he will make all the affairs belonging to the city-state subordinate to his own interests (Young, 2000). If such people are elected because of their wealth or power, not because of their virtue, we call this kind of government oligarchy.

On the one hand, Aristotle conducted in-depth research on the political system theories of important ancient Greek thinkers, including Socrates and Plato, attempting to find ideal political systems from their writings; On the other hand, he devoted great effort to conducting a detailed examination of the political systems of up to 158 city-states, attempting to find the optimal form of government from practice. His main purpose is to find the optimal form of government and prevent the worst political system (Young, 2000). It is precisely for this purpose that Aristotle used a large amount of space in his book “Politics” to repeatedly compare the advantages and disadvantages of various regimes, especially focusing on analyzing the best and worst regimes.

\section{Happy Life and Ideal Political System}

From Aristotle's extensive discourse on good and ideal regimes, we can summarize the three essential elements of the best regime that he diligently pursues. Firstly, “The best political system should be governed by the best people” (Robinson, 1971) The best political system must be ruled by the best people, and these rulers must be “talented and virtuous”, and everyone is convinced of their rule. Secondly, the best political system must have the best governance; In Aristotle's view, ‘the best regime' is almost directly equivalent to ‘the best governance'. Thirdly, and most importantly, the best political system must ensure that all citizens have a happy and beautiful life. Aristotle has always emphasized that the fundamental purpose of establishing a city-state is to enable people to live a happy and kind life. If a political system cannot maximize people's happiness, it cannot be considered a good political system at all. For Aristotle, an excellent government is the one that makes it possible for its citizens to attain the best life possible, and follows the best virtue.” (Robinson, 1971)

Aristotle did not seem to provide a clear answer to which of the three basic political systems he established, namely monarchy, aristocracy, and republic, was the best ideal political system. Moreover, his views often seem contradictory. He dismissed the idea of one person ruling because they would become a tyrant (Crisp, 2014). Later, he said that the highest power should be given to people with the highest talent and virtue. Among the three basic regimes, Aristotle found a monarch to be the best choice. At this time, he was clearly advocating that monarchy was the best form of government (Stewart, 1892) so monarchy and aristocracy are both regarded as the best regimes. Sometimes he said that republic was the right form of government, and democracy, which was ruled by citizens in turn, became the best form of government for the city-state. In addition, he said that the mixed government is the best regime. For example, he highly praised the mixed political system of Carthage.

Defining the “best” was a challenge for Aristotle. According to Voegelin (1987), in Aristotle's context at the time, the ‘best' did not only have one. meaning. It can refer to “the strongest,” “the healthiest,” “the most stable,” “the most suitable environment for achieving contemplative living,” “the most sufficient control over those who cannot achieve a good life,” “the maturity of the city-state,” or “the highest point in the realization of the essence of the city-state” (Voegelin, 1987). According to this understanding, what Aristotle called the best political system is actually the most suitable political system based on the specific conditions of each city state. As long as the regime adopted by the city-state is most suitable for promoting the overall public interest, this regime is the best system (Robinson, 1971). In other words, except for the three abnormal regimes that seek the personal interests of rulers, monarchy, aristocracy and republic can be the best regimes.

After a careful analysis of Aristotle's various statements about good government, we can still draw the conclusion that in Aristotle's mind, republican government or constitutional democracy is the best ideal political system. Aristotle did clearly say that monarchy and aristocracy were the best political systems, but in actual proof, he hardly took any monarchy and aristocracy that appeared in reality as examples of his good political systems. On the contrary, almost all the regimes he admires the most contain elements of republican democracy; The politicians he most admired were Solon and Pericles. He appreciates them mainly because they have carried out democratic reforms. After a detailed review of the history of the replacement of the constitutional system in Athens, Aristotle once again highly appreciated the democratic system in Athens. He praised the fact that people make themselves masters of everything in a democracy, using orders and courts to handle everything (Robinson, 1971).

Aristotle advocated republican democracy for the following four main reasons. Firstly, any good regime will eventually undergo mutation, and among all mutated regimes, the mutation of a republican regime is relatively least harmful. Secondly, under the conditions of republican democracy, the highest power of governance is exercised by citizens in rotation, which is most conducive to the political stability of the city-state. Thirdly, the power foundation of democratic politics is the middle class, and the power of the middle class is conducive to the long-term stability of the city-state (Robinson, 1971). Finally, extensive citizen participation is conducive to citizens taking responsibility for the city-state, thereby maximizing the public interests of the city-state (Robinson, 1971).

Most traditional views regard monarchy or aristocracy as Aristotle's ideal political system, while few people regard democracy as Aristotle's ideal political system. In contrast, contemporary Aristotelian researchers began to prefer democracy as Aristotle's ideal regime. For example, Michael Davis pointed out that Aristotle believed that all wise people's politics were ultimately unreliable, and political rule must obtain the consent of the governed, making democratic politics the ultimate ideal system (Dink, 1997).

Some critics of Aristotle’s views argue that he was wrong. For example, Moseley argues that Aristotle did not give trust to the monarch or any outstanding individual. He believed in the mean and the majority of citizens, and ultimately chose democratic politics. This view is somewhat surprising: although we might ultimately expect to understand his preference for hierarchical forms, this subtle shift—expressing democracy as the ideal—is unexpected (Moseley, 2007).

Josiah Ober also saw Aristotle's ambiguous attitude on the issue of the best regime, but Ober believed that democracy was at least one of the best regimes in Aristotle's mind. Ober said that Aristotle never specified which regime was the best, which is not difficult to understand. The best city-state is not only a natural body derived from human nature, but also a legislative design product based on the knowledge of Greek political history, which shows that it is both aristocracy and democracy. After comparing and analyzing the discourse on ideal political systems in the first and second volumes of Political Science, Ober pointed out that Aristotle regarded democracy as the best system in a certain analytical method, in the first volume's discourse on “naturally suitable” political systems and in the second volume's discourse on “the most ideal political system”. Because of its tolerance for citizens, democracy provides the answer to the problem of social stability, which is a prominent problem in Aristotle's Greek political history (Ober, 2015). Therefore, Aristotle gave a naturalized demonstration of the social, political and cultural advantages of democracy for social stability and long-term stability of the country in Politics.

If Aristotle's views on the best regime are still relatively vague and easily controversial, then his views on the worst regime are very clear and cannot tolerate any doubt. In his view, the worst regime in the world is the tyrant system or tyrant system as a pervert of monarchy. Tyrants not only monopolize power, dictatorship, and do whatever they want, but once they hold great power, they are not bound by any system and use their unlimited power to expand their desires, endangering the public interests of the city-state in various moving names. Throughout Aristotle's “Politics” and other related works, it is clear that he is full of hatred and contempt for tyrannical politics. Whenever there is an opportunity, he tirelessly reminds people that no one can hold unrestricted power under any name, otherwise it will cause endless harm to the city. He even publicly called on people to resist and kill the tyrant, giving the warriors who killed the tyrant supreme honor (Robinson, 1971).

In the final analysis, political systems are established and enforced by people, and the interaction between systems and people plays a crucial role in the political life of a city-state. Ignoring human factors, even with the best political system, cannot achieve the desired results. As Strauss said, although the system of government is an arrangement of political systems, but the political system reflects a more fundamental political reality - the relationship between authority and obedience among the different groups that make up the city-state (Young, 2000) In fact, the political system essentially reflects the relationship between the ruler and the ruled, and the political system is actually an institutional arrangement of the relationship between the ruler and the ruled. Aristotle deeply understood the dialectical relationship between system and human beings, and in addition to examining and analyzing political systems, he also expressed a series of important views on rulers and the ruled.

\section{Political System and Rulers}

The issue of which system and people are more important for the public interest of a city-state has been debated by thinkers before Aristotle. Some people believe that the legal system is more important, so they advocate that the highest power should belong to the law. Others believe that human factors are more important, and they advocate that the highest power should be handed over to the ruler. Aristotle said that no single persons should have all the power to rule others because they cannot avoid emotional impacts on their decisions. Therefore, it is better to entrust power to the law. However, the law itself may be biased towards the wealthy or towards the common people. To have a government based on laws that are biased towards the wealthy or towards the common people is no different from having the rule of an oligarchy or a democracy, in which the highest power is in the hands of either the few or the many” (Robinson, 1971). He clearly believes that both system and people are crucial, so he values both the role of the political system and the role of the ruler and the ruled.

One of Aristotle's most profound thoughts is to treat humans as political animals and see their inherent differences as political animals. However, distinguishing human beings as political animals from their nature and believing that humans are inherently unequal is Aristotle's greatest historical limitation. Aristotle divided humanity into three categories: men and women, free men and slaves, and civilized and barbarian people. Firstly, as animals, humans inevitably have a distinction between male and female, which is the difference between men and women. Due to limitations in physical and intellectual abilities, women rely on and are subject to men.

Due to significant group differences among humans as political animals, in order to maintain the basic order of social politics, it is naturally necessary to establish a basic relationship between ruling and being ruled among these groups. Aristotle believed that it is natural for men to rule over women, civilized men over barbarians, and especially free men over slaves. He repeatedly emphasized that humanity needs to live in a social and political community, and the basic order of each political community depends on the establishment of the basic relationship between the ruler and the ruled. A right regulation, however, of the relations between the two orders implies more than the mere fact of superiority on the part of the ruling element: there must also be a certain fitness and proportion between them. The people too must be fit for a share in the constitution; they must be both able and willing to perform their part in the common good (Young, 2000).

Aristotle regarded the distinction between rulers and ruled as an integral part of the natural order of the city-state, which was not only natural but also fair and reasonable. Under this premise, he further analyzed the basic relationship between the ruler and the ruled, and believed that there were three types of relationships between ruling and being ruled in a city-state, each of which was different. One is the ruling relationship between free people and slaves. In this relationship, the main purpose of rulers is for their own benefit. One type is the ruling relationship between father, wife, and children, in which the main purpose of the ruler is to serve the interests of the ruled. One type is the ruling relationship between citizens, in which the main purpose of rulers is to serve the public interest.

Regardless of the relationship between rule and being ruled mentioned above, the quality of the ruler should be extremely important. The quality of rulers is directly related to the safety of the political system and the interests of the rulers themselves. Aristotle believed that from the perspective of maintaining the public interest of the city-state government, the highest ruler of a city-state must possess three basic qualities. One is political quality, where rulers must be loyal to their own city-state; Secondly, excellent governance skills require rulers to be competent in the management of the city state; The third is noble moral qualities, and rulers must exemplary abide by the principles of fairness and justice (Robinson, 1971). These rules determine success.

Among all the qualities of rulers, Aristotle values their political and moral qualities the most. Aristotle believed that a good ruler should pursue a happy life for all citizens of the city-state, and use this as the fundamental criterion for judging the quality of rulers and their political systems. Therefore, in his view, the primary political quality of rulers is not to prioritize self-interest, but to prioritize the public interest of the city state, which is an unchangeable goal of good rulers (Robinson, 1971). Aristotle believed that there is a difference in the moral qualities between rulers and ruled individuals, and rulers should first possess the moral qualities of citizens. But this is far from enough. The moral character of rulers should surpass the level of ordinary people and become the most noble among citizens.

Although Aristotle emphasized the importance of rulers and hoped that there would be a person in the world who, since his virtue was superior to others and he was cautious and good, no one could surpass him, so that everyone could always follow and obey him, he was very clear that such a wise ruler did not actually exist in the real world. On the contrary, what he saw was the struggle for power and mutual strife between rulers. The endless power struggle between rulers is precisely the fundamental reason for the dynamic political situation and continuous internal strife in various city-states (Robinson, 1971) Therefore, Aristotle did not place all his political ideals on a “philosophical king” who combined wisdom and virtue, as his teacher Plato did. This may be an important reason why he admired republican democratic politics from the bottom of his bones.

Aristotle not only had no hope for any wise ‘philosophical king', but also had a special aversion to rulers who were obsessed with power and would not let go once they were in power. He compared the rulers who gained power by any means under the banner of the public interest of the city state, and compared the public power of the city state as their lifelong possession to “thieves and thugs”, and pointed out that defending the sins of these authoritarians is defending thieves and thugs. Since rulers are generally unreliable, in addition to seeking help from law and political system, Aristotle also sought help from citizens. Aristotle regarded the city-state as a political combination of citizens, who were the core roles in the political life of the city-state (Davis, 1996). Therefore, compared to scholars like Socrates and Plato, as well as other contemporary scholars, Aristotle particularly emphasized the role of citizens in the life of city-states and good governance.

As a great philosopher, Aristotle used the relationship between parts and the whole to describe the importance of citizens to a city state. He believes that a city-state is a political association of citizens, composed of specific citizens one by one. Citizens are a part of a city-state, while a city-state is the entirety of citizens. Just as the whole cannot leave its parts, without citizens, there is no city-state (Aristotle, Politics, Book III, Part 4). Aristotle not only defined the concept of citizen in political science, but also made clear distinctions between citizens and residents, citizens and expatriates. The citizens of a city-state are not ordinary residents, and the slaves within the city-state are not citizens. Citizenship is essentially a political concept, and its substantive element lies in the right to participate in the public political affairs of the city-state (Young, 2000).

Since the quality of citizens is directly related to the fate of a city-state, how to improve their quality naturally becomes extremely important. Aristotle believed that there are three factors that determine the quality of citizens, namely “external goods, bodily goods, and goods of the soul”. External goodness refers to material living conditions such as property; The “bodily goodness” refers to the physical and mental health of citizens; The ‘goodness of the soul' refers to the noble soul of a citizen, including its inherent qualities and talents.

Following the above logical thinking, Aristotle raised civic education to an unprecedented level. Citizen education is not only related to the physical and mental qualities of citizens themselves and their personal happiness, but also to the political fate and good life of the city-state. Even the rise, fall, and safety of the entire city-state are largely related to civic education, which has long been overlooked by people. Aristotle said: “The greatest of all the means towards ensuring stability in a constitution, and the one most neglected by most contemporary states, is to educate the citizens in the spirit of the constitution (the constitutional regime) (Davis, 1996). We must cultivate in the citizens those habits of conduct which will enable them to best preserve the form of government under which they live, whether it be a democracy or an oligarchy” (Robinson, 1971). Such habits must start with the leaders.

Aristotle believed that due to the different roles of rulers and ruled in the political life of the city-state, their requirements were also different, and therefore their education was not the same. However, as rulers are selected from among citizens, citizens should also possess certain knowledge of governance. In this way, civic education should include three aspects: firstly, general knowledge of city-state politics, including political and legal systems and administrative management; The second is the education of moral qualities, which enables citizens to develop the qualities of kindness, moderation, justice, courage, especially “obedience”; The third is ordinary national knowledge, which is a universal curriculum in school education, including four basic courses: reading and writing, gymnastics, music, and painting (Young, 2000). Aristotle placed the focus of civic education on young people, and among the existing 8 volumes of Political Science, there is a dedicated volume discussing “youth education”.

In other words, the political science of Aristotle that we are now reading is actually an unfinished work. His fundamental purpose in studying city-state regimes is to find which regime is the best, and what form is the best and possible system to achieve the ideal life that people imagine. What he referred to as the “good life” of a city-state is the “self-sufficient and supreme life” of the city-state, which is the “so-called true happiness and happiness of humanity”. Aristotle did not provide us with his final answer, and in fact, he could not provide us with the ultimate answer. Because on the one hand, this question is a dynamic one, with different answers for each era, and there is no ultimate universal answer; On the other hand, the answer to this question cannot be completed by any individual, and it requires the joint efforts of all humanity. In a sense, future thinkers and scholars continue to explore the best political system of humanity, striving to promote human political development towards democracy, peace, freedom, equality, fairness, justice, and good governance, which is precisely the continuation of a new chapter in Aristotle's Politics.

\begin{references}
\item Aristotle (1885). Politics (translated by Benjamin Jowett in 1999). \textit{Kitchener: Batoche Books}.
\item Crisp, R. (Ed.). (2014). \textit{Aristotle: Nicomachean Ethics}. Cambridge University Press.
\item Davis, M. (1996). \textit{The politics of philosophy: a commentary on Aristotle's Politics}. Rowman \& Littlefield.
\item Dink, M. (1997). The Politics of Philosophy: A Commentary on Aristotle's Politics.
\item Moseley, A. (2007). \textit{An introduction to political philosophy}. Bloomsbury Publishing.
\item Ober, J. (2015). Nature, history, and Aristotle's best possible regime. \textit{Aristotle’s ‘Politics’: A Critical Guide}.
\item Robinson, R. (1971). Aristotle's Politics, books iii and iv.
\item Stewart, J. A. (1892). \textit{Notes on the Nicomachean ethics of Aristotle} (Vol. 1). Clarendon Press.
\item Voegelin, E. (1987). \textit{The new science of politics: An introduction}. University of Chicago Press.
\item Young, C. M. (2000). Aristotle: Politics, books I and II. \textit{The Philosophical Review}, \textit{109}(1), 87-87.
\end{references}

\end{document}
